% !TEX program = xelatex
\documentclass[UTF8,a4paper,11pt]{ctexart}
\usepackage[a4paper,margin=2.0cm,headheight=14pt]{geometry}
\usepackage{amsmath,amssymb,booktabs,longtable,array,tabularx}
\usepackage{caption,graphicx,xcolor,xurl,hyperref}
\usepackage{url}
\definecolor{pyblue}{RGB}{36,91,163}
\hypersetup{colorlinks=true,linkcolor=pyblue,urlcolor=pyblue,
  pdftitle={PyQCU 与原生 QUDA pion 两点函数对照测试报告（P100 单/双卡阶梯）}}
\graphicspath{{../data/report_pion_c2_p100_20261014/}}
\newcommand{\code}[1]{\texttt{\detokenize{#1}}}
\setlength{\parindent}{2em}\setlength{\parskip}{0.35em}\emergencystretch=3em

\title{PyQCU(C++ MultiGrid) 与原生 QUDA 的 pion 两点函数对照测试报告\\\large P100 单卡/双卡阶梯与显存、时限边界}
\author{PyQCU 测试记录}
\date{2026 年 10 月 14 日}

\begin{document}
\maketitle

\begin{abstract}
本报告记录 \path{pyqcu/testing/quda/test_pion_c2.py} 在 NVIDIA P100（\code{sm_60}）上的单卡与双卡阶梯对照结果，
编排与产物格式沿用 \code{git tag test27}（commit \code{04382801c003353ecd1bc6b9bc10f6db750bf74d}）确立的口径：逐档保存求解总迭代数与纯求解耗时、
独立真残差、统一转为 QUDA/QDP 排布后的传播子差异，并记录峰值显存与容量/时限边界。
本次对照关闭 QUDA autotuning，两侧 MultiGrid 配置按粗层 color 对齐（QCU \code{dof_list=[12,48]}，
QUDA \code{n_vec=24}、\code{spin_block_size=2}，粗层 color 均为 48），并使用 \code{sm_60} 版
\code{libqcu.so} 与原生 \code{sm_60} QUDA 构建。结论：在 30 分钟/次的约束下，
\textbf{16$\times$16$\times$16$\times$32 是单卡与双卡均可完整跑通的最大档}；单卡在该档达到 97.9\% 显存占用，
双卡在 16$\times$16$\times$16$\times$48 档触及 30 分钟时限。
\end{abstract}

\section{环境与 provenance}
\label{sec:env}
\begin{table}[htbp]\centering\small
\caption{运行库与版本 provenance（完整 JSON：\texttt{provenance.json}）}
\begin{tabular}{@{}p{0.22\textwidth}p{0.70\textwidth}@{}}
\toprule
项 & 取值 \\ \midrule
QCU 库 & \path{data/build/libqcu_arch/sm60/libqcu.so}，SHA256 \code{e29c1cf04a57ec3b8af3208063456cab...}\\
 & SASS：\code{sm_60}（\code{cuobjdump --list-elf}）\\
QCU 运行时 & \code{torch 2.7.1+cu118}（\path{data/p100-site}，含 \code{sm_60} kernel image）+ \code{p100_env.sh} \\
QUDA & 原生 \code{QUDA 1.1.0 (git 1.1.0--sm_60)}，\code{pyquda 0.10.54}，SHA256 \code{aaf53934d5dab612e9bec02e20f20082...}\\
精度 & \code{complex64}（QUDA single），全链路一致 \\
计时口径 & \code{--quda-no-tuning}（默认关闭 autotuning）；warmup 与计时求解分离 \\
设备 & P100$\times$2（\code{CUDA_DEVICE_ORDER=PCI_BUS_ID}，device 0/1）\\
\bottomrule
\end{tabular}
\end{table}

\section{方法与公平性口径}
\label{sec:method}
\begin{itemize}
\item \textbf{物理量}：每个 \code{tsrc} 在空间原点投全 12 分量点源，求解 Clover-Wilson 系统
  $Dx=b$（$m=0.05$，$c_{sw}=1$，$\kappa=1/(2m+8)$），再按基脚本 GM-loop 结构做动量投影的
  pion 两点函数 $C_2(t)$（$\gamma_5\gamma_5$ 通道，零动量）。
\item \textbf{归一化}：两端算子满足 $D_{\mathrm{QUDA}}=(m+4)\,D_{\mathrm{PyQCU}}$；脚本不盲缩放，
  先用最小二乘拟合复数标度 $c$ 并报告 $c$ 与 $1/(2\kappa)=4.05$ 的相对偏差，再检查去掉 $c$ 后的残差。
\item \textbf{必需数值差异}：两侧传播子统一转换为 QUDA/QDP 排布
  $[\mathrm{row},\mathrm{parity},T,Z,Y,X/2,\mathrm{spin},\mathrm{color}]$ 后比较，并作为
  \code{qdp_layout_gate\{required:true\}} 硬闸（relL2 与 $\max|\Delta x|$ 双阈值），
  \code{--skip-c2} / \code{--skip-residual} 均不影响该项。
\item \textbf{MG 对齐}：粗化块 $[2,2,2,2]$、层数 2、粗层 color 48；两侧平滑器与迭代语义差异
  （QUDA 外层 GCR vs QCU 的 MG 重启迭代）在记录中显式标注，性能结论只采用同口径的纯求解耗时。
\item \textbf{多卡}：QCU 侧为单进程内“一线程一卡”并做跨卡一致性校验；QUDA 侧采用
  “一卡一进程 + \code{tsrc} 分片”，规避该 \code{sm_60} 构建中 MPI 集合调用的死锁（见第~\ref{sec:multi} 节）。
\end{itemize}

\section{逐档结果}
\label{sec:results}
\begin{table}[htbp]\centering\small
\caption{逐档求解规模与耗时（CSV：\texttt{units.csv}）。迭代数为该运行全部求解之和；
耗时只含纯求解，粗层构建/缓存加载单独记录在 \code{record.json}。}
\begin{tabular}{@{}llrrrrr@{}}
\toprule
格点 & 模式 & tsrc & QCU 迭代 & QCU s & QUDA 迭代 & QUDA s \\
\midrule
16$\times$16$\times$16$\times$32 & single & 1/1 & 48 & 0.1714159440016374 & 5 & 0.37051972502376884 \\
16$\times$16$\times$16$\times$32 & dual & 2/2 & 92 & 2.5351114311488345 & 10 & 0.44612990599125624 \\
16$\times$16$\times$16$\times$16 & single & 1/1 & 36 & 0.0566161279566586 & 6 & 0.07354366802610457 \\
16$\times$16$\times$16$\times$16 & single & 1/1 & 36 & 0.0566161279566586 & 6 & 0.07354366802610457 \\
16$\times$16$\times$16$\times$16 & single & 1/1 & 432 & 0.882596776005812 & 71 & 1.8289854740723968 \\
8$\times$8$\times$8$\times$16 & dual & 1/1 & 40 & 0.02875326806679368 & 6 & 0.060134319006465375 \\
8$\times$8$\times$8$\times$16 & single & 1/1 & 40 & 0.07552224001847208 & 6 & 0.021474626031704247 \\
\bottomrule
\end{tabular}
\end{table}

\begin{table}[htbp]\centering\small
\caption{逐档数值一致性：标度偏差、统一 QDP 排布的传播子差异与独立真残差。}
\begin{tabular}{@{}llrrrrr@{}}
\toprule
格点 & 模式 & $\delta c/c$ & QDP relL2 & QDP max$|\Delta x|$ & $r_{\mathrm{QCU}}$ & $r_{\mathrm{QUDA}}$ \\
\midrule
16$\times$16$\times$16$\times$32 & single & 1.5e-07 & 6.41e-06 & 2.39e-07 & 2.16e-07 & 4.95e-07 \\
16$\times$16$\times$16$\times$32 & dual & 1.4e-07 & 6.41e-06 & 4.79e-07 & 3.52e-07 & 4.95e-07 \\
16$\times$16$\times$16$\times$16 & single & 1.3e-08 & 2.87e-06 & 1.79e-07 & 2.62e-07 & 1.84e-07 \\
16$\times$16$\times$16$\times$16 & single & 1.3e-08 & 2.87e-06 & 1.79e-07 & 2.62e-07 & 1.84e-07 \\
16$\times$16$\times$16$\times$16 & single & 8.1e-08 & 4.67e-06 & 2.38e-07 & 2.34e-07 & 2.3e-07 \\
8$\times$8$\times$8$\times$16 & dual & 5.2e-07 & 7.6e-06 & 6.08e-07 & 2.13e-07 & 9.04e-07 \\
8$\times$8$\times$8$\times$16 & single & 4.7e-07 & 7.61e-06 & 6.17e-07 & 2.18e-07 & 8.95e-07 \\
\bottomrule
\end{tabular}
\end{table}

\noindent 另：16$^4$ 单卡 12 分量完整 pion $C_2$ 对照中，$C_2$ 相对差为 $2.100e-07$。

\begin{figure}[htbp]\centering
\includegraphics[width=0.92\textwidth]{fig_iterations.pdf}
\caption{逐档求解总迭代数（QUDA 为外层 GCR，QCU 为其 MG 重启/外层迭代，语义不同，仅作规模参考）。}
\end{figure}

\begin{figure}[htbp]\centering
\includegraphics[width=0.92\textwidth]{fig_solve_time.pdf}
\caption{纯求解耗时（对数轴，已排除 QUDA autotuning 与粗层构建）。}
\end{figure}

\begin{figure}[htbp]\centering
\includegraphics[width=0.92\textwidth]{fig_qdp_difference.pdf}
\caption{统一 QUDA/QDP 排布后的传播子差异与硬闸阈值（必需项）。}
\end{figure}

\begin{figure}[htbp]\centering
\includegraphics[width=0.92\textwidth]{fig_peak_memory.pdf}
\caption{PyQCU 阶段峰值显存（按设备取最大值）与 P100 16~GiB 上限对比。}
\end{figure}

\section{边界分析}
\label{sec:bounds}
\begin{enumerate}
\item \textbf{单卡 P100 显存边界}：$16\times16\times16\times32$ 档单卡运行峰值
  $16045/16384$~MiB（97.9\%），无 OOM；再放大 1.5 倍（$16^3\times48$）时粗层构建阶段
  已达 $14310$~MiB，单卡不可行。
\item \textbf{双卡 P100 时限边界}：$16^3\times48$（\code{mg_dof=[12,48]}）在
  \code{timeout 1800} 下 \code{exit=124}（30 分钟上限，无 OOM）。瓶颈为 QCU 粗层算子冷构建
  （批量化 stencil 探针，单卡完成），而非求解或显存。
\item \textbf{对齐与省时的取舍}：尝试 \code{mg_dof=[12,24]} 配合 QUDA \code{n_vec=12}
  （粗层 color 24）以减半构建时间，被 QUDA 拒绝：
  \code{Coarse Nc = 12 has not been instantiated (coarseColor=32)}，
  说明该构建不支持粗层 color 24；在不破坏 MG 对齐的前提下无法用降 E 换时间。
\end{enumerate}
边界明细：\path{data/report_pion_c2_p100_20261014/boundaries.csv}。

\section{多卡路径的修正过程}
\label{sec:multi}
QUDA 侧在该 \code{sm_60} 构建上，MPI 集合调用出现三类死锁（均由实测定位）：
(1) 非属主 rank 跳过 \code{invert} 导致属主 rank 在 MG 粗层通信中永久等待；
(2) rank 之间任务不对称时先到者执行 \code{MPI_Finalize} 而后者阻塞于 \code{Barrier}；
(3) 非属主 rank 跳过 \code{gather} 造成集合调用不对称。
最终采用\textbf{一卡一进程 + \code{tsrc} 分片}（\code{--quda-workers N}）实现多卡：
各进程单 rank、单卡、独立 \code{tsrc} 子集，产物按 worker 序号落盘后在对照阶段合并；
QCU 侧保持“一线程一卡”，两卡对同一 \code{tsrc} 复算的相对差为 $0.0$。
该改动同时修复了 worker 输出互相覆盖与旧产物混入两个数据完整性问题。

\section{复现命令}
\label{sec:repro}
\begin{verbatim}
source pyqcu/testing/qcu/strict/quda_comparison/p100_env.sh
export LD_LIBRARY_PATH=/root/PyQCU/data/build/libqcu_arch/sm60:$LD_LIBRARY_PATH
export CUDA_DEVICE_ORDER=PCI_BUS_ID

# 单卡（16^3x32，P100#0）
python pyqcu/testing/quda/test_pion_c2.py --lat 16 16 16 32 --tsrc 0 --components 1 \
  --qcu-devices 0 --qcu-threads 1 --quda-device 0 --skip-c2 --residual-rows 1 \
  --mg-nv-iters 1 --quda-warmup 0 \
  --mg-cache-dir data/testing/pion_c2/L16x16x16x32_dual/qcu/nullvec_cache --reuse-coarse \
  --run-id p100_single --run-dir data/testing/pion_c2/L16x16x16x32

# 双卡（16^3x32，两 worker + QCU 两线程）
python pyqcu/testing/quda/test_pion_c2.py --lat 16 16 16 32 --tsrc 0,1 --components 1 \
  --qcu-devices 0 1 --qcu-threads 2 --skip-c2 --residual-rows 1 --mg-nv-iters 1 --quda-warmup 0 \
  --multi-gpu --quda-workers 2 --quda-device 0 \
  --mg-cache-dir data/testing/pion_c2/L16x16x16x32_dual/qcu/nullvec_cache --reuse-coarse \
  --run-id L16x16x16x32_dual --run-dir data/testing/pion_c2/L16x16x16x32_dual
\end{verbatim}

\section{局限与后续}
\label{sec:limits}
\begin{itemize}
\item 迭代数语义不同（QUDA 外层 GCR / QCU MG 重启），跨后端只能作规模参考；性能结论采用同口径纯求解耗时。
\item 本轮只覆盖 \code{complex64}；\code{sm_60} QUDA 构建未启用 double-precision MultiGrid，
  c128 对照需改用 \code{QUDA_MULTIGRID_DOUBLE=ON} 的构建。
\item 继续放大需要把“粗层冷构建”与“对照”分离，或接入 dev80 的 HierarchicalCache
  （VRAM$\rightarrow$RAM$\rightarrow$DISK）；这两条路径尚未在本报告中验证。
\item 12 分量完整 pion $C_2$ 目前只在本机 16$^4$ 单卡档给出；$16^3\times32$ 档使用单分量诊断源，
  完整 12 分量在大格子上的对照需按时间预算另行安排。
\end{itemize}

\end{document}
